% Generated draft: requires mathematical and source review.
\documentclass[11pt]{article}
\usepackage[margin=25mm]{geometry}
\usepackage{fontspec}
\setmainfont{Latin Modern Roman}
\usepackage{amsmath,amssymb,mathtools,unicode-math}
\setmathfont{Latin Modern Math}
\usepackage{enumitem,needspace,xurl,hyperref}
\setlength{\parindent}{0pt}
\setlength{\parskip}{5pt}
\setlength{\emergencystretch}{3em}
\newcommand{\N}{\mathbb{N}}
\newcommand{\Z}{\mathbb{Z}}
\newcommand{\Q}{\mathbb{Q}}
\newcommand{\R}{\mathbb{R}}
\newcommand{\C}{\mathbb{C}}
\newcommand{\F}{\mathbb{F}}
\newcommand{\A}{\mathbb{A}}
\newcommand{\PP}{\mathbb{P}}
\newcommand{\E}{\mathbb{E}}
\newcommand{\eps}{\varepsilon}
\newcommand{\dd}{\,\mathrm{d}}
\begin{document}
\section*{UnsolvedMath ID: 801165}
\textbf{Source problem code:} catalogue entry
\subsection{Definitions and mathematical statement}
Use the standard dyadic convention: an orthonormal wavelet is a function \(\psi\in L^2(\R)\) for which
\[
\bigl\{2^{j/2}\psi(2^j x-k):j,k\in\Z\bigr\}
\]
is an orthonormal basis of \(L^2(\R)\). Use the Fourier transform \(\widehat f(\xi)=\int_{\R}f(x)e^{-2\pi i x\xi}\,\dd x\), extended to \(L^2(\R)\) by Plancherel's theorem. These conventions are consistent with the source's normalization of the smallest frequency-support measure as \(1\), but the supplied excerpt does not explicitly state the dilation or Fourier conventions.

A wavelet set is a Lebesgue measurable set \(W\subset\R\) of finite measure such that the function \(\psi_W\) defined by \(\widehat{\psi_W}=\mathbf 1_W\) is an orthonormal wavelet. In this normalization, such a set has Lebesgue measure \(1\).

\textbf{Problem (Larson, as reported by Bownik and Rzeszotnik, Problem 4).} Is it true that
\[
\forall\,\psi\in L^2(\R)\text{ that is an orthonormal wavelet},\quad
\exists\,W\subset\R\text{ that is a wavelet set},\quad
W\subset\operatorname{supp}\widehat\psi\,?
\]
The meaning of support and inclusion requires qualification because \(\widehat\psi\) is an \(L^2\) equivalence class. If support means the measurable nonzero set \(S_\psi=\{\xi:\widehat\psi(\xi)\ne0\}\), taken modulo null sets, the requested inclusion means
\[
\lambda(W\setminus S_\psi)=0,
\]
where \(\lambda\) denotes Lebesgue measure. If support instead means the closed essential support, it is the closed set
\[
C_\psi=\R\setminus\bigcup\{U\subset\R:U\text{ is open and }\widehat\psi=0\text{ almost everywhere on }U\},
\]
and the corresponding question requires \(W\subset C_\psi\), allowing modification of \(W\) on a null set. The supplied excerpt does not specify which support convention is intended; these readings are not silently identified here.

The source motivates the question by comparing two notions of minimal frequency support: attaining the smallest possible Lebesgue measure, stated there to be \(1\) and associated with minimally supported frequency (MSF) wavelets, and minimality under inclusion among frequency supports of orthonormal wavelets. It reports that whether these notions coincide was unknown in the source's account. This is supplied historical context, not an independently verified current assessment.
\subsection{Short English statement}
Larson's question asks whether the Fourier support of every orthonormal wavelet on \(\R\) contains a wavelet set. The draft uses the standard dyadic normalization and explicitly leaves the source's support convention unresolved. The registry labels the problem open; that label has not been independently verified.
\subsection{Sources}
\begin{itemize}
\item Local UnsolvedMath registry, permanent ID \(801165\), exact source problem code \texttt{catalogue entry}, title \emph{Open Problems in Wavelet Theory --- Problem 4}. The supplied record belongs to \emph{AMR Open Problem Lists} and carries the status \texttt{open}.
\item Bownik and Rzeszotnik, \emph{Open Problems in Wavelet Theory} (2019), Problem 4. Supplied source URL: \url{https://arxiv.org/abs/1905.08230}. The record reports access on 2026-10-16 and extraction by \texttt{arxiv\_latex}; these are supplied provenance details, not a claim of independent access or verification.
\item The statement attributes the problem to Larson, reportedly posed in the late 1990s and officially formulated in reference [La2]. The supplied excerpt does not provide the bibliographic entry for [La2]. The registry's proposer field names Bownik and Rzeszotnik and gives 2019, while its statement attributes the underlying question to Larson.
\end{itemize}
\end{document}
